\section{AMC 12 12A 2013}
        \begin{enumerate}
        	\item (\textit{AMC 12 12A 2013}) Square $ABCD$ has side length $10$. Point $E$ is on $\overline{BC}$, and the area of $\bigtriangleup ABE$ is $40$. What is $BE$?

                    \begin{figure}[H]
                        \centering
                        \includegraphics[width=0.5\textwidth]{Problem/AMC_12_12A_2013_1_Problem_Figure_1.png}
                    \end{figure}
                    
$\textbf{(A)} \ 4 \qquad \textbf{(B)} \ 5 \qquad \textbf{(C)} \ 6 \qquad \textbf{(D)} \ 7 \qquad \textbf{(E)} \ 8 \qquad$

 \textbf{Answer: }E



\textbf{Solution}

We are given that the area of $\triangle ABE$ is $40$, and that $AB = 10$.

 The area of a triangle is:

 $A = \frac{bh}{2}$

 Using $AB$ as the height of $\triangle ABE$,

 $40 = \frac{10b}{2}$

 and solving for $b$,

 $b = 8$, which is $\boxed{\textbf{(E)}}$



	\item (\textit{AMC 12 12A 2013}) A softball team played ten games, scoring $1,2,3,4,5,6,7,8,9$, and $10$ runs. They lost by one run in exactly five games. In each of the other games, they scored twice as many runs as their opponent. How many total runs did their opponents score? 


 $\textbf {(A) } 35 \qquad \textbf {(B) } 40 \qquad \textbf {(C) } 45 \qquad \textbf {(D) } 50 \qquad \textbf {(E) } 55$

 \textbf{Answer: }C



\textbf{Solution}

To score twice as many runs as their opponent, the softball team must have scored an even number.

 Therefore we can deduce that when they scored an odd number of runs, they lost by one, and when they scored an even number of runs, they won by twice as much.

 Therefore, the total runs by the opponent is $(2+4+6+8+10)+(1+2+3+4+5) = 45$, which is $C$



	\item (\textit{AMC 12 12A 2013}) A flower bouquet contains pink roses, red roses, pink carnations, and red carnations. One third of the pink flowers are roses, three fourths of the red flowers are carnations, and six tenths of the flowers are pink. What percent of the flowers are carnations?


 $\textbf{(A)}\ 15\qquad\textbf{(B)}\ 30\qquad\textbf{(C)}\ 40\qquad\textbf{(D)}\ 60\qquad\textbf{(E)}\ 70$

 \textbf{Answer: }E



\textbf{Solution}

We are given that $\frac{6}{10} = \frac{3}{5}$ of the flowers are pink, so we know $\frac{2}{5}$ of the flowers are red.

 Since $\frac{1}{3}$ of the pink flowers are roses, $\frac{2}{3}$ of the pink flowers are carnations.

 We are given that $\frac{3}{4}$ of the red flowers are carnations.

 The percent of flowers that are carnations is 

 $\frac{3}{5} * \frac{2}{3} + \frac{2}{5} * \frac{3}{4} = \frac{2}{5} + \frac{3}{10} = \frac{7}{10} = 70\%$, which is $E$



	\item (\textit{AMC 12 12A 2013}) What is the value of 
 \[\frac{2^{2014}+2^{2012}}{2^{2014}-2^{2012}}?\]
$\textbf{(A)}\ -1\qquad\textbf{(B)}\ 1\qquad\textbf{(C)}\ \frac{5}{3}\qquad\textbf{(D)}\ 2013\qquad\textbf{(E)}\ 2^{4024}$

 \textbf{Answer: }C



\textbf{Solution}

$\frac{2^{2014}+2^{2012}}{2^{2014}-2^{2012}}$

 We can factor a ${2^{2012}}$ out of the numerator and denominator to obtain

 $\frac{2^{2012}*(2^2+1)}{2^{2012}*(2^2-1)}$

 The ${2^{2012}}$ cancels, so we get 

 $\frac{(2^2+1)}{(2^2-1)}=\frac{5}{3}$, which is $C$




\textbf{Solution 2}

Suppose that 
 \[x=2^{2012}.\] Then the given expression is equal to 
 \[\frac{4x+x}{4x-x}=\frac{5x}{3x}=\boxed{\frac{5}{3}\textbf{(C)}}.\]




\textbf{Solution 3 (Elimination)}

Choice $A$ cannot be true, because $2^{2014}+2^{2012}$ is clearly larger than $2^{2014}-2^{2012}$. We can apply our same logic to choice $B$ and eliminate it as well. $2013$ and $2^{4024}$ are all whole numbers, but $2^{2014}+2^{2012}$ is not a multiple of $2^{2014}-2^{2012}$, so we can eliminate choices $D$ and $E$ too. This leaves us with choice $\boxed{C}$ as our final answer.

 


	\item (\textit{AMC 12 12A 2013}) Tom, Dorothy, and Sammy went on a vacation and agreed to split the costs evenly. During their trip Tom paid $$105$, Dorothy paid $$125$, and Sammy paid $$175$. In order to share the costs equally, Tom gave Sammy $t$ dollars, and Dorothy gave Sammy $d$ dollars. What is $t-d$?


 $\textbf{(A)}\ 15\qquad\textbf{(B)}\ 20\qquad\textbf{(C)}\ 25\qquad\textbf{(D)}\ 30\qquad\textbf{(E)}\ 35$

 \textbf{Answer: }B



\textbf{Solution 1}

Simply write down two algebraic equations. We know that Tom gave $t$ dollars and Dorothy gave $d$ dollars. In addition, Tom originally paid $105$ dollars and Dorothy paid $125$ dollars originally. Since they all pay the same amount, we have: 
 \[105 + t = 125 + d.\] Rearranging, we have 
 \[t-d = \boxed{\textbf{(B)} 20}.\]

 Solution $\textcopyright 2018$ RandomPieKevin. All Rights Reserved. skrrt...

 



\textbf{Solution 2}

Add up the amounts that Tom, Dorothy, and Sammy paid to get \$$405$, and divide by 3 to get \$$135$, the amount that each should have paid.

 Tom, having paid \$$105$, owes Sammy \$$30$, and Dorothy, having paid \$$125$, owes Sammy \$$10$.

 Thus, $t - d = 30 - 10 = 20$, which is $\boxed{\textbf{(B)}}$



	\item (\textit{AMC 12 12A 2013}) In a recent basketball game, Shenille attempted only three-point shots and two-point shots. She was successful on $20\%$ of her three-point shots and $30\%$ of her two-point shots. Shenille attempted $30$ shots. How many points did she score?


 $\textbf{(A)}\ 12\qquad\textbf{(B)}\ 18\qquad\textbf{(C)}\ 24\qquad\textbf{(D)}\ 30\qquad\textbf{(E)}\ 36$

 \textbf{Answer: }B



\textbf{Solution}

Let the number of 3-point shots attempted be $x$. Since she attempted 30 shots, the number of 2-point shots attempted must be $30 - x$.

 Since she was successful on $20\%$, or $\frac{1}{5}$, of her 3-pointers, and $30\%$, or $\frac{3}{10}$, of her 2-pointers, then her score must be

 $\frac{1}{5}*3x + \frac{3}{10}*2(30-x)$
$\frac{3}{5}*x + \frac{3}{5}(30-x)$
$\frac{3}{5}(x+30-x)$
$\frac{3}{5}*30$
$18$, which is $B$




\textbf{Solution 2}

Since the problem doesn't specify the number of 3-point shots she attempted, it can be assumed that number doesn't matter, so let it be $0$. Then, she must have attempted $30$ 2-point shots. So, her score must be:

 $\frac{3}{10}*30*2$,which is $B$.

 



\textbf{Solution 3}

(similar to Solution 1, however a slightly more obvious way)

 Say that 

 x = \# of 2-pt shots

 y = \# of 3-pt shots

 Because the total number of shots is $30$, $x + y = 30$

 However, Shenille was only successful on $20\%$ of the 3-pt shots, and $30\%$ of the 2-pt shots, so 

 $0.2x + 0.3y$ = \#number of successful shots

 For each successful shot, there is an associated number of points with it. 

 Therefore, $0.2(x)(3) + (0.3)(y)(2)$ = her score

 this evaluates to $0.6 (x + y)$ = her score

 $x + y$ is already determined to be 30, so her score is $0.6 (30) = \boxed{\textbf{(B)}\ 18}$

 \textbf{Additional note} It is also reasonably easy to find all possibilities for the number of two-point and three-point shots she made. Just note that both numbers of successful throws have to be integers. For "$30\%$ of her two-point shots" to be an integer we need the number of two-point shots to be divisible by 10. This only leaves four possibilities for the number of two-point shots: 0, 10, 20, or 30. Each of them also works for the three-point shots, and as shown above, for each of them the total number of points scored is the same.

 


	\item (\textit{AMC 12 12A 2013}) The sequence $S_1, S_2, S_3, \cdots, S_{10}$ has the property that every term beginning with the third is the sum of the previous two.  That is, 
 \[S_n = S_{n-2} + S_{n-1} \text{ for } n \ge 3.\] Suppose that $S_9 = 110$ and $S_7 = 42$.  What is $S_4$?


 $\textbf{(A)}\ 4\qquad\textbf{(B)}\ 6\qquad\textbf{(C)}\ 10\qquad\textbf{(D)}\ 12\qquad\textbf{(E)}\ 16\qquad$

 \textbf{Answer: }C



\textbf{Solutions 1 and 2}

Solution 1:

 $S_9 = 110$, $S_7 = 42$
$S_8 = S_9 - S_ 7 = 110 - 42 = 68$
$S_6 = S_8 - S_7 = 68 - 42 = 26$
$S_5 = S_7 - S_6 = 42 - 26 = 16$
$S_4 = S_6 - S_5 = 26 - 16 = 10$

 Therefore, the answer is $\boxed{\textbf{(C) }{10}}$

 
== Solution 2 ==(first time writing something like this, sorry if bad structure)(no I did not write solution 1)

 S3=S1+S2, S4=S2+S1+S2=2*S1 +S2, S5=3*S1+2*S2, S6=5*S1+3*S2, S7=8*S1+5*S2, S8=13*S1+8*S2, S9=21*S1+13*S2. S9=110, S7=42. 8*S1+5*S2=42. 21*S1+13*S2=110. 105*S1+65*S2=550, 104*S1+65*S2=546. S1=550-546=4. S2:21*4+13*S2=110. 13*S2=26. S2=2. S3=S1+S2=6. S4=S2+S3= $\boxed{\textbf{(C) }{10}}$ (idk how to open a new section so sorry) 

 


	\item (\textit{AMC 12 12A 2013}) Given that $x$ and $y$ are distinct nonzero real numbers such that $x+\frac{2}{x} = y + \frac{2}{y}$, what is $xy$?


 $\textbf{(A)}\ \frac{1}{4}\qquad\textbf{(B)}\ \frac{1}{2}\qquad\textbf{(C)}\ 1\qquad\textbf{(D)}\ 2\qquad\textbf{(E)}\ 4\qquad$

 \textbf{Answer: }D



\textbf{Solution 1}

$x+\frac{2}{x}= y+\frac{2}{y}$

 Since $x\not=y$, we may assume that $x=\frac{2}{y}$ and/or, equivalently, $y=\frac{2}{x}$.

 Cross multiply in either equation, giving us $xy=2$.

 $\boxed{\textbf{(D) }{2}}$




\textbf{Solution 2}

$x+\frac{2}{x}= y+\frac{2}{y}$
$x-y+\frac{2}{x}-\frac{2}{y} = 0$
$(x-y)+2(\frac{y-x}{xy}) = 0$
$(x-y)(1-\frac{2}{xy})=0$

 Since $x\not=y$
$1 = \frac{2}{xy}$
$xy = 2$




\textbf{Solution 3}

Let $A = x + \frac{2}{x} = y + \frac{2}{y}.$ Consider the equation 
 \[u + \frac{2}{u} = A.\] Reorganizing, we see that $u$ satisfies 
 \[u^2 - Au + 2 = 0.\] Notice that there can be at most two distinct values of $u$ which satisfy this equation, and $x$ and $y$ are two distinct possible values for $u.$ Therefore, $x$ and $y$ are roots of this quadratic, and by Vieta’s formulas we see that $xy$ thereby must equal $\boxed{2}.$




\textbf{Solutions 4 and 5}


 \[x + \frac{2}{x} = y + \frac{2}{y}.\]

 Multiply both sides by xy to get

 
 \[x^2y + 2y = y^2x +2x\]

 Rearrange to get

 
 \[x^2y - y^2x = 2x - 2y\]

 Factor out $xy$ on the left side and $2$ on the right side to get

 
 \[xy(x-y) = 2(x-y)\]

 Divide by $(x-y)$ {You can do this since x and y are distinct} to get

 $\boxed{\textbf{(D) }{2}}$

 Solution 5:(sorry I don’t know how to open a new section whoever sees this next pls help!)(This is not a rigorous solution but it is extremely quick)(no I am not the person who wrote solution 4)

 Question states x+2/x =y+2/y. Simply let x equal 1, y equal two. It satisfies because 3=3. Since it didn’t specify any further, we may realize immediately that xy=1*2= $\boxed{\textbf{(D) }{2}}$



	\item (\textit{AMC 12 12A 2013}) In $\triangle ABC$, $AB=AC=28$ and $BC=20$.  Points $D,E,$ and $F$ are on sides $\overline{AB}$, $\overline{BC}$, and $\overline{AC}$, respectively, such that $\overline{DE}$ and $\overline{EF}$ are parallel to $\overline{AC}$ and $\overline{AB}$, respectively.  What is the perimeter of parallelogram $ADEF$?


 
                    \begin{figure}[H]
                        \centering
                        \includegraphics[width=0.5\textwidth]{Problem/AMC_12_12A_2013_9_Problem_Figure_1.png}
                    \end{figure}
                    
$\textbf{(A) }48\qquad \textbf{(B) }52\qquad \textbf{(C) }56\qquad \textbf{(D) }60\qquad \textbf{(E) }72\qquad$

 \textbf{Answer: }C



\textbf{Solution}

Note that because $\overline{DE}$ and $\overline{EF}$ are parallel to the sides of $\triangle ABC$, the internal triangles $\triangle BDE$ and $\triangle EFC$ are similar to $\triangle ABC$, and are therefore also isosceles triangles.

 It follows that $BD = DE$. Thus, $AD + DE = AD + DB = AB = 28$.

 Since opposite sides of parallelograms are equal, the perimeter is $2 * (AD + DE) =   \boxed{\textbf{(C) }{56}}$.

 



\textbf{Solution (Cheese)}

We can set point $F$ to be on point $C$, and point $D$ to be on point $A$.

 This makes a degenerate parallelogram with sides of length $28$ and $0$, so it has a perimeter of $28 + 28 =  \boxed{\textbf{(C) }{56}}$.

 


	\item (\textit{AMC 12 12A 2013}) Let $S$ be the set of positive integers $n$ for which $\frac{1}{n}$ has the repeating decimal representation $0.\overline{ab} = 0.ababab\cdots,$ with $a$ and $b$ different digits.  What is the sum of the elements of $S$?


 $\textbf{(A)}\ 11\qquad\textbf{(B)}\ 44\qquad\textbf{(C)}\ 110\qquad\textbf{(D)}\ 143\qquad\textbf{(E)}\ 155\qquad$

 \textbf{Answer: }D



\textbf{Solution 1}

Note that $\frac{1}{11} = 0.\overline{09}$.

 Dividing by 3 gives $\frac{1}{33} = 0.\overline{03}$, and dividing by 9 gives $\frac{1}{99} = 0.\overline{01}$.

 $S = \{11, 33, 99\}$
$11 + 33 + 99 = 143$

 The answer must be at least $143$, but cannot be $155$ since no $n \le 12$ other than $11$ satisfies the conditions, so the answer is $143$.

 



\textbf{Solution 2}

Let us begin by working with the condition $0.\overline{ab} = 0.ababab\cdots,$. Let $x = 0.ababab\cdots$. So, $100x-x = ab \Rightarrow x = \frac{ab}{99}$. In order for this fraction $x$ to be in the form $\frac{1}{n}$, $99$ must be a multiple of $ab$. Hence the possibilities of $ab$ are $1,3,9,11,33,99$. Checking each of these, $\frac{1}{99} = 0.\overline{01}, \frac{3}{99}=\frac{1}{33} = 0.\overline{03}, \frac{9}{99}=\frac{1}{11} = 0.\overline{09}, \frac{11}{99}=\frac{1}{9} = 0.\overline{1}, \frac{33}{99} =\frac{1}{3}= 0.\overline{3},$ and $\frac{99}{99} = 1$. So the only values of $n$ that have distinct $a$ and $b$ are $11,33,$ and $99$. So, $11+33+99= \boxed{\textbf{(D)} 143}$




\textbf{Solution 3}

Notice that we have $\frac{100}{n}= ab.\overline{ab}$ 

 We can subtract $\frac{1}{n}=00.\overline{ab}$ to get 
 \[\frac{99}{n}=ab\]

 From this we determine $n$ must be a positive factor of $99$

 
The factors of $99$ are $1,3,9,11,33,$ and $99$.

 For $n=1,3,$ and $9$ however, they yield $ab=99,33$ and $11$ which doesn't satisfy $a$ and $b$ being distinct. 

 For $n=11,33$ and $99$ we have $ab=09,03$ and $01$.  (Notice that $a$ or $b$ can be zero)

 The sum of these $n$ are $11+33+99=143$
$\boxed{\textbf{(D)} 143}$




\textbf{Solution 4}

As in previous solutions, we have $n|99$ and $\overline{ab} = 99/n$. If we had $a=b$, the decimal would be $0.\overline{a}$, which is characterized by $n|9$ and $a = 9/n$. So we seek the sum of the factors of 99 that are not also factors of 9.

 Since $99 = 3^2 \cdot 11$, the sum is $(1 + 3 + 9)(1 + 11) - (1 + 3 + 9) = 13(12 - 1) = \textbf{(D)} 143$.

 


	\item (\textit{AMC 12 12A 2013}) Triangle $ABC$ is equilateral with $AB=1$. Points $E$ and $G$ are on $\overline{AC}$ and points $D$ and $F$ are on $\overline{AB}$ such that both $\overline{DE}$ and $\overline{FG}$ are parallel to $\overline{BC}$. Furthermore, triangle $ADE$ and trapezoids $DFGE$ and $FBCG$ all have the same perimeter. What is $DE+FG$?


 
                    \begin{figure}[H]
                        \centering
                        \includegraphics[width=0.5\textwidth]{Problem/AMC_12_12A_2013_11_Problem_Figure_1.png}
                    \end{figure}
                    
$\textbf{(A) }1\qquad \textbf{(B) }\dfrac{3}{2}\qquad \textbf{(C) }\dfrac{21}{13}\qquad \textbf{(D) }\dfrac{13}{8}\qquad \textbf{(E) }\dfrac{5}{3}\qquad$

 \textbf{Answer: }C



\textbf{Solution}

Let $AD = x$, and $AG = y$. We want to find $DE + FG$, which is nothing but $x+y$.

 Based on the fact that $ADE$, $DEFG$, and $BCFG$ have the same perimeters, we can say the following:

 $3x = x + 2(y-x) + y = y + 2(1-y) + 1$

 Simplifying, we can find that

 $3x = 3y-x = 3-y$

 Since $3-y = 3x$, $y = 3-3x$.

 After substitution, we find that $9-10x = 3x$, and $x$ = $\frac{9}{13}$.

 Again substituting, we find $y$ = $\frac{12}{13}$.

 Therefore, $x+y$ = $\frac{21}{13}$, which is $C$



	\item (\textit{AMC 12 12A 2013}) The angles in a particular triangle are in arithmetic progression, and the side lengths are $4,5,x$. The sum of the possible values of $x$ equals $a+\sqrt{b}+\sqrt{c}$ where $a, b$, and $c$ are positive integers. What is $a+b+c$?


 $\textbf{(A)}\ 36\qquad\textbf{(B)}\ 38\qquad\textbf{(C)}\ 40\qquad\textbf{(D)}\ 42\qquad\textbf{(E)}\ 44$

 \textbf{Answer: }A



\textbf{Solution}

Because the angles are in an arithmetic progression, and the angles add up to $180^{\circ}$, the second largest angle in the triangle must be $60^{\circ}$. Also, the side opposite of that angle must be the second longest because of the angle-side relationship. Any of the three sides, $4$, $5$, or $x$, could be the second longest side of the triangle.

 The law of cosines can be applied to solve for $x$ in all three cases.

 When the second longest side is $5$, we get that $5^2 = 4^2 + x^2 - 2(4)(x)\cos 60^{\circ}$, therefore $x^2 - 4x - 9 = 0$. By using the quadratic formula, $x = \frac {4 + \sqrt{16 + 36}}{2}$, therefore $x = 2 + \sqrt{13}$.

 When the second longest side is $x$, we get that $x^2 = 5^2 + 4^2 - 40\cos 60^{\circ}$, therefore $x = \sqrt{21}$.

 When the second longest side is $4$, we get that $4^2 = 5^2 + x^2 - 2(5)(x)\cos 60^{\circ}$, therefore $x^2 - 5x + 9 = 0$. Using the quadratic formula, $x = \frac {5 + \sqrt{25 - 36}}{2}$. However, $\sqrt{-11}$ is not real, therefore the second longest side cannot equal $4$.

 Adding the two other possibilities gets $2 + \sqrt{13} + \sqrt{21}$, with $a = 2, b=13$, and $c=21$. $a + b + c = 36$, which is answer choice $\text{(A)}$.

 



\textbf{Solution 2 (Similar to [1]) and Solution 3(non trig solution)}

The angle measures of the triangle are $\theta$, $\theta + d$, and $\theta + 2d$ for some common difference $d$ through the definition of arithmetic sequence.

 Angle sum formula gives $3\theta + 3d = 180$ and thus $\theta+d = 60$.

 One inscribed measure of the triangle is $\theta + d = 60^\circ$. We do casework on where this angle is inscribed, using LoC along the way. Note that $\cos(60^\circ) = \frac{1}{2}$

 If $60^\circ$ is inscribed with 4 and 5, then $x^2 = 41 - 20 = 21 \implies x = \sqrt{21}$.

 If $60^\circ$ is inscribed with 5 and $x$, then $-9 = x^2 - 5x$ and thus $x^2 - 5x +9 = 0$. No such solutions exist checking the discriminant.

 If $60^\circ$ is inscribed with 4 and $x$, then $9 = x^2 - 4x$ and thus $x^2 - 4x - 9 = 0$. It follows that $x_{+} = 2+\sqrt{13}$ is the only correct root that satisfies the properties of Plane Euclidean Geometry. 

 
 \[\sum \textbf{all~} x = 2+\sqrt{13} + \sqrt{21} \implies 2 + 13 + 21 = \boxed{36}\]

 Solution 3:(solution 2 was not written by me, whoever sees this next please fix the horrible structure! Idk how to open a new page)(also I don’t know how to LaTeX diagrams! HELP)

 First, set the angles equal to a-b, a, a+b. Added together should equal 180. So, 3a=180, a=60. So the angles are 60,60-b, 60+b for some number b. We already know two of the sides are four and five, so now there are three cases. 1. Four and five are the two sides around the 60 degrees, 2.only the 4 side is touching the 60 degrees, 3. Only 5 side is touching the 60 degrees.

 Case 1: both 4 side and 5 side touch 60 degrees. Now, draw an altitude from the vertice the 4 side does not share with the five side. It will result in a 30-60-90 triangle with 4 as the hypotenuse, so the part on the. 5 side equals 2 and the height of the triangle, both the 30-60-90 and the whole triangle, is 2sqrt(3). We have a right triangle from the altitude, the remaining 3 on the 5 side, and the remaining side of the big triangle that we are looking for. By Pythagorean theorem, (2sqrt(3))\textasciicircum{}2+3\textasciicircum{}2=21=remaining side squared. So in this case the third side is equal to sqrt(21).

 Case 2: only 4 side is touching the 60 degrees. Set third side we are looking for as x. Once again draw an altitude down, this time from the shared vertice of the 4 and 5 sides to side x. Once again height equals 2sqrt(3) with an almost exact copy paste of above. Once again the bottom side part is 2 for the 30-60-90 triangle. However this time, the Pythagorean theorem equation we get is (2sqrt(3))\textasciicircum{}2 + (a-2)\textasciicircum{}2 =5\textasciicircum{}2? This means that (a-2)\textasciicircum{}2 equals 13. Thus, a\textasciicircum{}2-4a-9=0. Use quadratic formula to get that the roots are 2+sqrt(13),

 \begin{verbatim}
2-sqrt(13). Clearly a length must be positive so we discard 2-sqrt(13). So in case 2, the third side is equal to 2-sqrt(13).
\end{verbatim}

Case 3: set unknown third side as a. Once again draw an altitude down, this time from the connecting vertice of an and the 4 side to the 5 side. The bottom part on the 5 side equals a/2, and the height of the triangles is (sqrt(3)/2)*a. So the Pythagorean equation we receive this time is ((sqrt(3)/2)*a)\textasciicircum{}2 +(5-a/2)\textasciicircum{}2 =16. expanding it out results in (3/4)*a\textasciicircum{}2 + 25-5a+(a\textasciicircum{}2/4) =16. Which, simplified, equals a\textasciicircum{}2-5a+9=0. However, the discriminant is -11 which is clearly less than zero so there are no real roots. If you draw this case out you will realize it is impossible to draw such a triangle. 

 Thus our answer is 2+sqrt(13)+sqrt(21)=a+sqrt(b)+sqrt(c) which means a=2, b=13, c=21. So, a+b+c equals 2+13+21=36 which is answer choice A. 

 


	\item (\textit{AMC 12 12A 2013}) Let points $A = (0,0) , \ B = (1,2), \ C = (3,3),$ and $D = (4,0)$. Quadrilateral $ABCD$ is cut into equal area pieces by a line passing through $A$. This line intersects $\overline{CD}$ at point $\left (\frac{p}{q}, \frac{r}{s} \right )$, where these fractions are in lowest terms. What is $p + q + r + s$?


 $\textbf{(A)} \ 54 \qquad \textbf{(B)} \ 58 \qquad  \textbf{(C)} \ 62 \qquad \textbf{(D)} \ 70 \qquad \textbf{(E)} \ 75$

 \textbf{Answer: }B



\textbf{Solution}

If you have graph paper, use Pick's Theorem to quickly and efficiently find the area of the quadrilateral. If not, just find the area by other methods.

 Pick's Theorem states that

 $A$ = $I$ $+$ $\frac{B}{2}$ - $1$, where $I$ is the number of lattice points in the interior of the polygon, and $B$ is the number of lattice points on the boundary of the polygon.

 In this case, 

 $A$ = $5$ $+$ $\frac{7}{2}$ - $1$ = $7.5$

 so

 $\frac{A}{2}$ = $3.75$

 The bottom half of the quadrilateral makes a triangle with base $4$ and half the total area, so we can deduce that the height of the triangle must be $\frac{15}{8}$ in order for its area to be $3.75$. This height is the y coordinate of our desired intersection point.

 
Note that segment CD lies on the line $y = -3x + 12$. Substituting in $\frac{15}{8}$ for y, we can find that the x coordinate of our intersection point is $\frac{27}{8}$.

 Therefore the point of intersection is ($\frac{27}{8}$, $\frac{15}{8}$), and our desired result is $27+8+15+8= \boxed{ \text{B}) \ 58}$.

 



\textbf{Solution 2 (Shoelace)}

Let the point of intersection be $E$, with coordinates $(x, y)$. Then, $ABCD$ is cut into $ABCE$ and $AED$. 

 Since the areas are equal, we can use \href{https://artofproblemsolving.com/wiki/index.php?title=Shoelace\_Theorem}{Shoelace Theorem} to find the area. This gives $3 + 3x - 3y = 4y$.

 The line going through $CD$ is $y = -3x + 12$. Since $E$ is on $CD$, we can substitute this in, giving $3 + 3x = -21x + 84$. Solving for $x$ gives $\frac{27}{8}$. Plugging this back into the line equation gives $y = \frac{15}{8}$, for a final answer of $\boxed{58}$.

 


	\item (\textit{AMC 12 12A 2013}) The sequence


 $\log_{12}{162}$, $\log_{12}{x}$, $\log_{12}{y}$, $\log_{12}{z}$, $\log_{12}{1250}$


 is an arithmetic progression. What is $x$?


 $\textbf{(A)} \ 125\sqrt{3} \qquad \textbf{(B)} \ 270 \qquad \textbf{(C)} \ 162\sqrt{5} \qquad \textbf{(D)} \ 434 \qquad \textbf{(E)} \ 225\sqrt{6}$

 \textbf{Answer: }B



\textbf{Solution 1}

Since the sequence is arithmetic, 

 $\log_{12}{162}$ + $4d$ = $\log_{12}{1250}$, where $d$ is the common difference.

 
Therefore,

 $4d$ = $\log_{12}{1250}$ - $\log_{12}{162}$ = $\log_{12}{(1250/162)}$, and

 $d$ = $\frac{1}{4}$($\log_{12}{(1250/162)}$) = $\log_{12}{(1250/162)^{1/4}}$

 


 Now that we found $d$, we just add it to the first term to find $x$:

 $\log_{12}{162}$ + $\log_{12}{(1250/162)^{1/4}}$ = $\log_{12}{((162)(1250/162)^{1/4})}$
$x$ = $(162)$$(1250/162)^{1/4}$ = $(162)$$(625/81)^{1/4}$ = $(162)(5/3)$ = $270$, which is $\boxed{B}$




\textbf{Solution 2}

As the sequence $\log_{12}{162}$, $\log_{12}{x}$, $\log_{12}{y}$, $\log_{12}{z}$, $\log_{12}{1250}$ is an arithmetic progression, the sequence $162,x,y,z,1250$ must be a geometric progression.

 If we factor the two known terms we get $162=2\cdot 3^4$ and $1250=2\cdot 5^4$, thus the quotient is obviously $5/3$ and therefore $x=162\cdot(5/3) = 270$.

 


	\item (\textit{AMC 12 12A 2013}) Rabbits Peter and Pauline have three offspring—Flopsie, Mopsie, and Cotton-tail. These five rabbits are to be distributed to four different pet stores so that no store gets both a parent and a child. It is not required that every store gets a rabbit. In how many different ways can this be done?


 $\textbf{(A)} \ 96 \qquad  \textbf{(B)} \ 108 \qquad  \textbf{(C)} \ 156 \qquad  \textbf{(D)} \ 204 \qquad  \textbf{(E)} \ 372$

 \textbf{Answer: }D



\textbf{Solution 1}

There are two possibilities regarding the parents.

 1) Both are in the same store. In this case, we can treat them both as a single bunny, and they can go in any of the 4 stores. The 3 baby bunnies can go in any of the remaining 3 stores. There are $4 \cdot 3^3 = 108$ combinations.

 
2) The two are in different stores. In this case, one can go in any of the 4 stores, and the other can go in any of the 3 remaining stores. The 3 baby bunnies can each go in any of the remaining 2 stores. There are $4 \cdot 3 \cdot 2^3 = 96$ combinations.

 Adding up, we get $108 + 96 = 204$ combinations.

 



\textbf{Solution 2}

We tackle the problem by sorting it by how many stores are involved in the transaction.

 1) $2$ stores are involved. There are $\binom{4}{2} = 6$ ways to choose which of the stores are involved and 2 ways to choose which store recieves the parents. $6 \cdot 2 = 12$ total arrangements.

 
2) $3$ stores are involved. There are $\binom{4}{3} = 4$ ways to choose which of the stores are involved. We then break the problem down to into two subsections - when the parents and grouped together or sold separately.

 
\textbf{Separately:}All children must be in one store. There are $3!$ ways to arrange this. $6$ ways in total.

 \textbf{Together:}Both parents are in one store and the 3 children are split between the other two. There are $\binom{3}{2}$ ways to split the children and $3!$ ways to choose to which store each group will be sold. $3! \cdot \binom{3}{2} = 18$.

 $(6 + 18) \cdot 4 = 96$ total arrangements.

 
3) All $4$ stores are involved. We break down the problem as previously shown.

 
\textbf{Separately:}All children must be split between two stores. There are $\binom{3}{2} = 3$ ways to arrange this. We can then arrange which group is sold to which store in $4!$ ways. $4! \cdot 3 = 72$.

 \textbf{Together:}Both parents are in one store and the 3 children are each in another store. There are $4! = 24$ ways to arrange this.

 $24 + 72 = 96$ total arrangements.

 
\textbf{Final Answer:}$12 + 96 + 96 = \boxed{\textbf{(D)} \: 204}$.

 


	\item (\textit{AMC 12 12A 2013}) $A$, $B$, $C$ are three piles of rocks. The mean weight of the rocks in $A$ is $40$ pounds, the mean weight of the rocks in $B$ is $50$ pounds, the mean weight of the rocks in the combined piles $A$ and $B$ is $43$ pounds, and the mean weight of the rocks in the combined piles $A$ and $C$ is $44$ pounds. What is the greatest possible integer value for the mean in pounds of the rocks in the combined piles $B$ and $C$?


 $\textbf{(A)} \ 55 \qquad \textbf{(B)} \ 56 \qquad \textbf{(C)} \ 57 \qquad \textbf{(D)} \ 58 \qquad \textbf{(E)} \ 59$

 \textbf{Answer: }E



\textbf{Solution 1}

Let pile $A$ have $A$ rocks, and so on.

 The total weight of $A$ and $C$ can be expressed as $44(A + C)$.

 To get the total weight of $B$ and $C$, we add the weight of $B$ and subtract the weight of $A$: $44(A + C) + 50B - 40A = 4A + 44C + 50B$

 Therefore, the mean of $B$ and $C$ is $\frac{4A + 44C + 50B}{B + C}$, which is simplified to $44 + \frac{4A + 6B}{B + C}$.

 We now need to eliminate $A$ in the numerator.  Since we know that $40A + 50B = 43(A + B)$, we have $A = \frac{7}{3}B$

 Substituting,

 $44 + \frac{4(\frac{7}{3}B) + 6B}{B + C}=44 + \frac{46}{3}*\frac{B}{B + C}$

 In order to maximize $\frac{B}{B + C}$, we can minimize the denominator by letting $C = 1$ (C must be a positive integer). Since $\frac{46}{3}$ must cancel to give an integer, and the only fraction that satisfies both conditions is $\frac{45}{46}$

 Plugging in, we get

 $44 + \frac{46}{3} * \frac{45}{46} = 44 + 15 = 59$, which is choice E

 



\textbf{Solution 2}

Suppose there are $A,B,C$ rocks in the three piles, and that the mean of pile C is $x$, and that the mean of the combination of $B$ and $C$ is $y$. We are going to maximize $y$, subject to the following conditions:

 
 \[40A+50B=43(A+B)\]
 \[40A+xC=44(A+C)\]
 \[50B+xC=y(B+C)\]

 which can be rearranged as:

 
 \[7B=3A\]
 \[(x-44)C=4A\]
 \[(x-y)C=(y-50)B\]

 Let us test $y=59$ is possible. If so, it is already the answer. If not, there will be some contradiction. So the third equation becomes 

 
 \[(x-59)C=9B.\]

 So $15C = (x-44)C - (x-59)C = 4A - 9B$, $45C=4(3A)-27B=28B-27B$, $105C=28A-9(7B)=A$, therefore,

 $A=105C, B=45C, x=4(105)+44=464$, which gives us a consistent solution. Therefore $y=59$ is the answer.

 (Note: To further illustrate the idea, let us look at $y=60$ and see what happens. We then get $7\cdot 16C = 4A-30A<0$, which is a contradiction!)

 



\textbf{Solution 3}

Obtain the 3 equations as in \textbf{solution 2}.

 
 \[7B=3A\]
 \[(x-44)C=4A\]
 \[(x-y)C=(y-50)B\]

 Our goal is to try to isolate $y$ into an inequality. The first equation gives $A=\frac{7}{3}B$, which we plug into the second equation to get

 
 \[(x-44)C=\frac{28}{3}B\]

 To eliminate $x$, subtract equation 3 from equation 2:

 
 \[(x-44)C-(x-y)C=\frac{28}{3}B-(y-50)B\]
 \[(y-44)C=(\frac{178}{3}-y)B\]

 In order for the coefficients to be positive, 
 \[44<y<\frac{178}{3}\]

 Thus, the greatest integer value is $y=59$, choice $(E)$.

 



\textbf{Solution 4}

Let the number of rocks in $A$ be $a$, $B$ be $b$, $C$ be $c$. The total weight of $A$ be $40a$, $B$ be $50b$, $C$ be $kc$.

 We can write the information given as, $\frac{40a + 50b}{a+b} = 43$, $\frac{40a + kc}{a+c} = 44$, $\frac{50b + kc}{b+c} = ?$
$40a + 50b = 43 a + 43 b$, $3a = 7b$
$40a + kc = 44a + 44 c$, $kc = 4a + 44c = \frac{28}{3}b + 44c$
$\frac{50b + kc}{b+c} = \frac{50 \cdot \frac{3}{7}a + kc}{\frac{3}{7}a+c} = \frac{150a + 7kc}{3a + 7c} = \frac{150a + 28a + 308c}{3a+7c} = \frac{178a + 308c}{3a+7c} = \frac{44(3a+7c)+46a}{3a+7c}$
$= 44 + \frac{46a}{3a+7c} = 44 + \frac{46}{3 + \frac{7}{a}} < 44 + \frac{46}{3} \approx 59.3$
$\frac{50b + kc}{b+c} \le \boxed{\textbf{(E) } 59}$




\textbf{Solution 5}

Let the total number of rocks in pile $A$ be $A_n$, and the total number of rocks in pile $B$ be $B_n$. Then, by restriction 3 (the average of $A$ and $B$), we can establish the equation: 
 \[\frac{40A_n+50B_n}{A_n+B_n}=43\].Cross-multiplying, we get: 
 \[40A_n+50B_n=43A_n+43B_n \implies 3A_n=7B_n\].Let's say we have $7k$ rocks in $A$ and $3k$ rocks in $B$. Hence, we have $280k$ and $150k$ as the total weight of piles $A$ and $B$, respectively. Let the total weight of $C$ be $m$, and the total number of rocks in $C$ be $n$. Using the last restriction regarding the average of piles $A$ and $C$, we have: 
 \[\frac{280k+m}{7k+n}=44 \implies 280k + m=280k + 28k + 44n \implies m=28k+44n\].To find the average of piles $B$ and $C$, we can establish the expression: 
 \[\frac{150k+28k+44n}{3k+n}=\frac{178k+44n}{3k+n}=\frac{132k+44n+46k}{3k+n}=\frac{44(3k+n)+46k}{3k+n}=44+\frac{46k}{3k+n}.\]When we let the final expression equal $59$, we get: 
 \[44+\frac{46k}{3k+n}=59\implies\frac{46k}{3k+n}=15\]Cross-multiplying, we get: 
 \[46k=45k+14n\implies k=14n\]$k$ is still positive here, so 59 works. As this is the greatest option, we can circle $\textbf{(E)}$ immediately.To show why $59$ is the greatest, consider the following:When we let the final expression equal $60$, we get: 
 \[44+\frac{46k}{3k+n}=60\implies\frac{46k}{3k+n}=16\]Cross-multiplying, we get: 
 \[46k=48k+14n\implies k=-7n\]Since $k$ is positive, the final expression could not equal 60. It further implies that the final expression could not equal any other integer greater than 60. Therefore, we have our final answer $\boxed{59}$.

 



\textbf{Solution 6}

Denote the number of rocks in each of the three piles as $a$, $b$, and $c$, respectively, and denote the weight of each rock in pile $C$ as $x$. We are given:\begin{align}\frac{40a + 50b}{a+b} &= 43\\\frac{40a + xc}{b+c} &= 44\\\end{align}From these equations, we get $b = \frac{3}{7}a$ and $c = \frac{4}{x-44}a$.

 Now, we can compute the average weight in combined piles B and C:\begin{align}\frac{50b + xc}{b+c} &= \frac{\frac{150}{7}a + \frac{4x}{x-44}a}{\frac{3}{7}a + \frac{4}{x-44}a}\\&= \frac{178x-6600}{3x-104}\\&= 59 + \frac{1}{3} - \frac{C}{3x-104}\\\end{align}Where $C$ is just some constant that I'm too lazy to compute. As $x$ approaches infinity, we have that $\frac{C}{3x-104}$ becomes negligibly small, and at some point, $59$ is attained. ~akliu

 


	\item (\textit{AMC 12 12A 2013}) A group of $12$ pirates agree to divide a treasure chest of gold coins among themselves as follows. The $k^\text{th}$ pirate to take a share takes $\frac{k}{12}$ of the coins that remain in the chest. The number of coins initially in the chest is the smallest number for which this arrangement will allow each pirate to receive a positive whole number of coins. How many coins does the $12^{\text{th}}$ pirate receive?


 $\textbf{(A)} \ 720 \qquad  \textbf{(B)} \ 1296 \qquad  \textbf{(C)} \ 1728 \qquad  \textbf{(D)} \ 1925 \qquad  \textbf{(E)} \ 3850$

 \textbf{Answer: }D



\textbf{Solution 1}

The first pirate takes $\frac{1}{12}$ of the $x$ coins, leaving $\frac{11}{12} x$.

 The second pirate takes $\frac{2}{12}$ of the remaining coins, leaving $\frac{10}{12}\cdot   \frac{11}{12}*x$.

 Note that

 $12^{11} = (2^2 \cdot 3)^{11} = 2^{22} \cdot 3^{11}$
$11! = 11 \cdot 10 \cdot 9 \cdot 8 \cdot 7 \cdot 6 \cdot 5 \cdot 4 \cdot 3 \cdot 2$

 
All the $2$s and $3$s cancel out of $11!$, leaving

 $11 \cdot 5 \cdot 7 \cdot 5 = 1925$

 in the numerator.

 
We know there were just enough coins to cancel out the denominator in the fraction. So, at minimum, $x$ is the denominator, leaving $\boxed{\textbf{(D) }1925}$ coins for the twelfth pirate.

 



\textbf{Solution 2}

The answer cannot be an even number. Here is why:

 Consider the highest power of 2 that divides the starting number of coins, and consider how this value changes as each pirate takes their share. At each step, the size of the pile is multiplied by some $\frac{n}{12}$. This means the highest power of 2 that divides the number of coins is continually decreasing or staying the same (except once, briefly, when we multiply by $\frac{2}{3}$ for pirate 4, but it immediately drops again in the next step).

 Therefore, if the 12th pirate's coin total were even, then it can't be the smallest possible value, because we can safely cut the initial pot (and all the intermediate totals) in half. We could continue halving the result until the 12th pirate's total is finally odd.

 Only one of the choices given is odd, $\boxed{\textbf{(D) }1925}$.

 Note that if there had been multiple odd answer choices, we could repeat the reasoning above to see that the value also must not be divisible by 3.

 



\textbf{Solution 3 (Work backwards)}

Let $x$ be the number of coins the $12$th pirate takes. Then the number of coins the $k<12$th pirate takes is $\frac{12}{1} \cdot \frac{12}{2} \cdots \frac{12}{12-k} x$. For all these to be an integer, we need the denominators to divide into the numerators. Looking at prime factors, obviously there are sufficient $2$s and $3$s, so we just need $5^2 \cdot 7 \cdot 11 = \boxed{\textbf{(D) }1925}$ to divide into $x$. -Frestho

 


	\item (\textit{AMC 12 12A 2013}) Six spheres of radius $1$ are positioned so that their centers are at the vertices of a regular hexagon of side length $2$. The six spheres are internally tangent to a larger sphere whose center is the center of the hexagon. An eighth sphere is externally tangent to the six smaller spheres and internally tangent to the larger sphere. What is the radius of this eighth sphere?


 $\textbf{(A)} \ \sqrt{2} \qquad \textbf{(B)} \ \frac{3}{2} \qquad \textbf{(C)} \ \frac{5}{3} \qquad \textbf{(D)} \ \sqrt{3} \qquad \textbf{(E)} \ 2$

 \textbf{Answer: }B



\textbf{Solution 1}

Set up an isosceles triangle between the center of the 8th sphere and two opposite ends of the hexagon. Then set up another triangle between the point of tangency of the 7th and 8th spheres, and the points of tangency between the 7th sphere and 2 of the original spheres on opposite sides of the hexagon. Express each side length of the triangles in terms of r (the radius of sphere 8) and h (the height of the first triangle). You can then use Pythagorean Theorem to set up two equations for the two triangles, and find the values of h and r.

 $(1+r)^2=2^2+h^2$
$(3\sqrt{2})^2=3^2+(h+r)^2$
$r = \boxed{\textbf{(B) }\frac{3}{2}}$




\textbf{Solution 2}

We have a regular hexagon with side lengths 2 and six spheres on each vertex with radius 1 that are internally tangent, therefore drawing radii going through all of them would create this regular hexagon.

 There is a larger sphere which the 6 spheres are internally tangent to, with center in the center of the hexagon. To find the radius of the larger sphere we must first, either by prior knowledge or by deducing from the angle sum that the hexagon can be split into 6 equilateral triangles from it's vertices, that the radius is $2+1=3$

 The 8th sphere is now, when thinking about it in 3D, sitting on top of the 6 spheres, which is the only possibility for it to tangent all the 6 small spheres externally and the larger sphere internally. The ring of the 6 small spheres is symmetrical and the 8th sphere will be resting with it's center aligned with the diameter of the large sphere.

 We can therefore now create a triangle with the horizontal component 2, as it is from the vertex of the hexagon to the center of the hexagon. The vertical component is from the center of the large sphere to the center of the 8th sphere. This length equals 3, the radius of the large sphere, take away the radius of the 8th sphere, we can call it r, since the radius of the large sphere will include the diameter of the 8th sphere if we subtract radius we will reach the center.The last component is the hypotenuse of the right angled triangle. This consists of the radius of the small sphere - 1 - and the radius of the 8th sphere - r -.

 We therefore now have a right angled triangle which when applied Pythagoras states $2^2+(3-r)^2=(1+r)^2$Expanding brackets gives us $4+9-6r+r^2=1+2r+r^2$ here we can cancel out $r^2$Isolating the r's $12=8r$and then finally we have the answer: $r=\frac{12}{8}=\frac{3}{2}$




\textbf{Solution 3 (Simply Drawing)}


                    \begin{figure}[H]
                        \centering
                        \includegraphics[width=0.5\textwidth]{Solution/AMC_12_12A_2013_18_Solution_Figure_1.png}
                    \end{figure}
                    It can be seen that the diameter of the eighth sphere is equal to the radius of the seventh sphere by drawing out a diagram of the insides of the seventh sphere. The radius of the seventh sphere is $2+1=3$, the radius of the eight sphere is $\boxed{\textbf{(B) } \frac32}$.

 


	\item (\textit{AMC 12 12A 2013}) In $\bigtriangleup ABC$, $AB = 86$, and $AC = 97$. A circle with center $A$ and radius $AB$ intersects $\overline{BC}$ at points $B$ and $X$. Moreover $\overline{BX}$ and $\overline{CX}$ have integer lengths. What is $BC$?


 $\textbf{(A)} \ 11 \qquad  \textbf{(B)} \ 28 \qquad  \textbf{(C)} \ 33 \qquad  \textbf{(D)} \ 61 \qquad  \textbf{(E)} \ 72$

 \textbf{Answer: }D



\textbf{Solution 1 (Diophantine PoP)}


                    \begin{figure}[H]
                        \centering
                        \includegraphics[width=0.5\textwidth]{Solution/AMC_12_12A_2013_19_Solution_Figure_1.png}
                    \end{figure}
                    

 Let circle $A$ intersect $AC$ at $D$ and $E$ as shown. We apply Power of a Point on point $C$ with respect to circle $A.$ This yields the diophantine equation 

 
 \[CX \cdot CB = CD \cdot CE\]
 \[CX(CX+XB) = (97-86)(97+86)\]
 \[CX(CX+XB) = 3 \cdot 11 \cdot 61.\]

 Since lengths cannot be negative, we must have $CX+XB \ge CX.$ This generates the four solution pairs for $(CX,CX+XB)$: 
 \[(1,2013) \qquad (3,671) \qquad (11,183) \qquad (33,61).\]

 However, by the Triangle Inequality on $\triangle ACX,$ we see that $CX>13.$ This implies that we must have $CX+XB= \boxed{\textbf{(D) }61}.$ 

 (Solution by unknown, latex/asy modified majorly by samrocksnature)

 



\textbf{Solution 2}

Let $BX = q$, $CX = p$, and $AC$ meet the circle at $Y$ and $Z$, with $Y$ on $AC$.  Then $AZ = AY = 86$.  Using the Power of a Point, we get that $p(p+q) = 11(183) = 11 * 3 * 61$.  We know that $p+q>p$, and that $p>13$ by the triangle inequality on $\triangle ACX$.  Thus, we get that $BC = p+q = \boxed{\textbf{(D) }61}$




\textbf{Solution 3}

Let $x$ represent $CX$, and let $y$ represent $BX$. Since the circle goes through $B$ and $X$, $AB = AX = 86$.Then by Stewart's Theorem,

 $xy(x+y) + 86^2 (x+y) = 97^2 y + 86^2 x.$
$x^2 y + xy^2 + 86^2 x + 86^2 y = 97^2 y + 86^2 x$
$x^2 + xy + 86^2 = 97^2$

 (Since $y$ cannot be equal to $0$, dividing both sides of the equation by $y$ is allowed.)

 $x(x+y) = (97+86)(97-86)$
$x(x+y) = 2013$

 The prime factors of $2013$ are $3$, $11$, and $61$. Obviously, $x < x+y$. In addition, by the Triangle Inequality, $BC < AB + AC$, so $x+y < 183$. Therefore, $x$ must equal $33$, and $x+y$ must equal $\boxed{\textbf{(D) }61}$




\textbf{Solution 4}

Motivation and general line of reasoning: we use a law of cosines condition on triangle $ABX$ and triangle $ABC$ to derive some equivalent formations of the same quantity $\cos B$, which looks promising since it involves the desired length $BC$, as well as $BX$ and $CX$.An intermediate step would be to use the integer condition and pay attention to the divisors of stuff. 

 Let $BX$=$x_1$ and $CX$=$x_2$. 

 First we have $\cos(B)=\frac{86^2+x_{1}^2-86^2}{172x_{1}}$ by applying the law of cosines to triangle $ABX$. Another equivalent formation of $\cos B$ is$\frac{86^2+(x_{1}+x_2)^2-97^2}{172(x_{1}+x_2)}$. Now that we have the necessary ingredients, we can make a system of equations and deduce that $x_1=\frac{(x_1+x_2)^2+86^2-97^2}{x_1+x_2}$. 

 Don't lose focus by now-we try to find $x_1+x_2$, an integer value as given in the problem. To do this, we want the quantity $\frac{(x_1+x_2)^2+86^2-97^2}{x_1+x_2}$ to 1) be an integer and 2) smaller than $x_1+x_2$. For the sake of conciseness in notation we let $M=x_1+x_2$, then $M+\frac{86^2-97^2}{M}$ is an integer. Now recalling the fact that $(a+b)(a-b)=a^2-b^2$, we get that $\frac{183(-11)}{M}$ must be an integer.

 Now the prime factor decomposition of $183 \cdot -11$ is  $(61)(3)(-11)$. Trying out all the possible integer values that divide this quantity, we get that the only viable option for $M$ is 61 (verify that yourself!) Therefore the answer is $\boxed{\textbf{(D) }61}$. (Solution by CreamyCream123)

 


	\item (\textit{AMC 12 12A 2013}) Let $S$ be the set $\{1,2,3,...,19\}$. For $a,b \in S$, define $a \succ b$ to mean that either $0 < a - b \le 9$ or $b - a > 9$. How many ordered triples $(x,y,z)$ of elements of $S$ have the property that $x \succ y$, $y \succ z$, and $z \succ x$?


 $\textbf{(A)} \ 810 \qquad  \textbf{(B)} \ 855 \qquad  \textbf{(C)} \ 900 \qquad  \textbf{(D)} \ 950 \qquad  \textbf{(E)} \ 988$

 \textbf{Answer: }B



\textbf{Solution 1}

Imagine that the 19 numbers are just 19 persons sitting evenly around a circle $C$; each of them is facing to the center. 

 One may check that $x \succ y$ if and only if $y$ is one of the 9 persons on the left of $x$, and $y \succ x$ if and only if $y$ is one of the 9 persons on the right of $x$. Therefore, "$x \succ y$ and $y \succ z$ and $z \succ x$" implies that $x, y, z$ cuts the circumference of $C$ into three arcs, each of which has no more than $10$ numbers sitting on it (inclusive).

 We count the complement: where the cut generated by $(x,y,z)$ has ONE arc that has more than $10$ persons sitting on. Note that there can only be one such arc because there are only $19$ persons in total.

 Suppose the number of persons on the longest arc is $k>10$. Then two places of $x,y,z$ are just chosen from the two end-points of the arc, and there are $19-k$ possible places for the third person. Once the three places of $x,y,z$ are chosen, there are three possible ways to put $x,y,z$ into them clockwise. Also, note that for any $k>10$, there are $19$ ways to choose an arc of length $k$. Therefore the total number of ways (of the complement) is 

 
 \[\sum_{k=11}^{18} 3\cdot 19 \cdot (19-k) = 3\cdot 19 \cdot (1+\cdots+8) = 3\cdot 19\cdot 36\]

 So the answer is

 
 \[3\cdot \binom{19}{3} - 3\cdot 19\cdot 36 = 3\cdot 19 \cdot (51 - 36) = 855\]

 NOTE: this multiple-choice problem can be done even faster -- after we realized the fact that each choice of the three places of $x,y,z$ corresponds to $3$ possible ways to put them in, and that each arc of length $k>10$ has $19$ equitable positions, it is evident that the answer should be divisible by $3\cdot 19$,  which can only be $855$ from the five choices.

 



\textbf{Solution 2}

First, we can find out that the only $x, y, z$ that satisfy the conditions in the problem are $(1 \leq z-y \leq 9 , 1 \leq y-x \leq 9 , z-x \geq 10)$ , $(1 \leq y-x \leq 9 , 1 \leq x-z \leq 9 , y-z \geq 10)$ , and $(1 \leq x-z \leq 9 , 1 \leq z-y \leq 9 , x-y \geq 10)$. 

 Consider the 1st set of conditions for $x, y, z$. We get that there are

 
 \[45\cdot 10 - \sum_{k=1}^{9} \binom{k+1}{2} = 285\]

 cases for the first set of conditions.

 Since the 2nd and 3rd set of conditions are simply rotations of the 1st set, the total number of cases is

 
 \[3\cdot 285 = \fbox{(B)855}\]




\textbf{Solution 3 (Common sense)}

The conditions are very strange. The major way $x,y,z$ can satisfy is if $(1 \leq z-y \leq 9 , 1 \leq y-x \leq 9 , z-x \geq 10)$ because at this point $z > x$ no matter what. This is the only case that works, but we can rotate the variables around, so our answer is a multiple of $3$.

 Suppose $z=y+a$ and $y=x+b$ where $0<a,b<10$. That means $z=x+a+b$ and $10 \leq z \leq 18$. WLOG $z>y>x$. Start casework on what $a+b$ is, anywhere from $10$ to $18$. There are 2 main factors:1. the number of ways to choose $a,b$, which is $a+b-1$.2. the number of ways to choose $x$ or start the sequence, which is $19-(a+b)$The sum of all cases turns out to be the sum of squares from $9^2$ to $1^2$, which is $\frac{9 \cdot 10 \cdot 19}{2 \cdot 3}$.Multiply that by $3$ and get $\fbox{(B)855}$




\textbf{Solution 4 (Modular Arithmetic)}

Note that if we take $a, b, c\mod{19}$, the conditions are equivalent to $0<a-b\le{9}\mod{19}$, $0<b-c\le{9}\mod{19}$, and $0<c-a\le{9}\mod{19}$. Let $m=a-b, n=b-c$. We have $19$ ways to pick $a$, and $81$ cases in our sample space for $m, n$. Of these cases, only those where $m+n>9$ will work. By stars and bars, this turns out to be $\binom{10}{2}=45$, hence the answer is $45*19=\fbox{(B)855}$.~sigmapie

 



\textbf{Solution 5 (pretty braindead)}

$x \succ_1 y$ will denote the first inequality, and $x\succ_2 y$ will denote the second.

 
We soon see that we cannot have $x\succ_1 y \succ_1 z \succ_1 x$ by adding the inequalities we obtain from them together, as it results in $0<0$.

 
Similarly, we can't have $x\succ_2 y \succ_2 z \succ_2 x$. $x\succ_2 y \succ_2 z \succ_1 x$ is also not possible by observing the inequalities.

 
Thus, we only have to consider $x \succ_1 y\succ_1 z \succ_2 x$. Something like $x \succ_2 y \succ_1 z \succ_1 x$ is just a rotation of this case.

 
Here, we start with $0<x-y\le 9$, $0<y-z\le 9$, and $x-z > 9$. We begin with casework on $z$, since it is the smallest.

 
When $z=1$, $x$ can range from $[11,19]$. There are $9+8+7+6+5+4+3+2+1 = 45$ different values for $y$ as x ranges across $11-19$.

 
When $z=2$, $x$ can range from $[12-19]$. There are $9+8+7+6+5+4+3+2 = 44$ different values for $y$ as x ranges across $12-19$.

 
Continuing this pattern, we get the sum $45+44+42+39+35+30+24+17+9 = 285$. Multiplying by three to count for rotations, we get $\boxed{(\textbf{B}) 855}$.

 -skibbysiggy

 


	\item (\textit{AMC 12 12A 2013}) Consider $A = \log (2013 + \log (2012 + \log (2011 + \log (\cdots + \log (3 + \log 2) \cdots ))))$. Which of the following intervals contains $A$?


 $\textbf{(A)} \ (\log 2016, \log 2017)$
$\textbf{(B)} \ (\log 2017, \log 2018)$
$\textbf{(C)} \ (\log 2018, \log 2019)$
$\textbf{(D)} \ (\log 2019, \log 2020)$
$\textbf{(E)} \ (\log 2020, \log 2021)$

 \textbf{Answer: }A



\textbf{Solution 1}

Let $f(x) = \log(x + f(x-1))$ and $f(2) = \log(2)$, and from the problem description, $A = f(2013)$

 We can reason out an approximation, by ignoring the $f(x-1)$:

 $f_{0}(x) \approx \log x$

 And a better approximation, by plugging in our first approximation for $f(x-1)$ in our original definition for $f(x)$:

 $f_{1}(x) \approx \log(x + \log(x-1))$

 And an even better approximation:

 $f_{2}(x) \approx \log(x + \log(x-1 + \log(x-2)))$

 Continuing this pattern, obviously, will eventually terminate at $f_{x-1}(x)$, in other words our original definition of $f(x)$.

 However, at $x = 2013$, going further than $f_{1}(x)$ will not distinguish between our answer choices. $\log(2012 + \log(2011))$ is nearly indistinguishable from $\log(2012)$.

 So we take $f_{1}(x)$ and plug in.

 $f(2013) \approx \log(2013 + \log 2012)$

 Since $1000 < 2012 < 10000$, we know $3 < \log(2012) < 4$. This gives us our answer range:

 $\log 2016 < \log(2013 + \log 2012) < \log(2017)$
$(\log 2016, \log 2017)$




\textbf{Solution 2}

Suppose $A=\log(x)$. Then $\log(2012+ \cdots)=x-2013$. So if $x>2017$, then $\log(2012+\log(2011+\cdots))>4$. So $2012+\log(2011+\cdots)>10000$. Repeating, we then get $2011+\log(2010+\cdots)>10^{7988}$. This is clearly absurd (the RHS continues to grow more than exponentially for each iteration).So, $x$ is not greater than $2017$. So $A<\log(2017)$. But this leaves only one answer, so we are done.

 



\textbf{Solution 3}

Define $f(2) = \log(2)$, and $f(n) = \log(n+f(n-1)), \text{ for } n > 2.$ We are looking for $f(2013)$. First we show

 \textbf{Lemma.} For any integer $n>2$, if $n < 10^k-k$ then $f(n) < k$.

 \textbf{Proof.} First note that $f(2) < 1$. Let $n<10^k-k$. Then $n+k<10^k$, so $\log(n+k)< k$. Suppose the claim is true for $n-1$. Then $f(n) = \log(n+f(n-1)) < \log(n + k) < k$.  The Lemma is thus proved by induction.

 Finally, note that $2012 < 10^4 - 4$ so that the Lemma implies that $f(2012) < 4$. This means that $f(2013) = \log(2013+f(2012)) < \log(2017)$, which leaves us with only one option $\boxed{\textbf{(A) } (\log 2016, \log 2017)}$.

 



\textbf{Solution 4}

Define $f(2) = \log(2)$, and $f(n) = \log(n+f(n-1)), \text{ for } n > 2.$ We start with a simple observation:

 \textbf{Lemma.} For $x,y>2$, $\log(x+\log(y)) < \log(x)+\log(y)=\log(xy)$.

 \textbf{Proof.} Since $x,y>2$, we have $xy-x-y = (x-1)(y-1) - 1 > 0$, so $xy - x-\log(y) > xy - x - y > 0$.

 It follows that $\log(z+\log(x+\log(y))) < \log(x)+\log(y)+\log(z)$, and so on. 

 Thus $f(2010) < \log 2 + \log 3 + \cdots + \log 2010 < \log 2010 + \cdots + \log 2010 <  2009\cdot 4 = 8036$. 

 Then $f(2011) = \log(2011+f(2010)) < \log(10047) < 5$. 

 It follows that $f(2012) = \log(2012+f(2011)) < \log(2017) < 4$. 

 Finally, we get $f(2013) = \log(2013 + f(2012)) < \log(2017)$, which leaves us with only option $\boxed{\textbf{(A)}}$.

 



\textbf{Solution 5 (nonrigorous + abusing answer choices.)}

Intuitively, you can notice that $\log(2012+\log(2011+\cdots(\log(2)\cdots))<\log(2013+\log(2012+\cdots(\log(2))\cdots))$, therefore (by the answer choices) $\log(2012+\log(2011+\cdots(\log(2)\cdots))<\log(2021)$. We can then say:

 
 \[x=\log(2013+\log(2012+\cdots(\log(2))\cdots))\]
 \[\log(2013)<x<\log(2013+\log(2021))\]
 \[\log(2013)<x<\log(2013+4)\]
 \[\log(2013)<x<\log(2017)\]

 The only answer choice that is possible given this information is $\boxed{\textbf{(A)}}$




\textbf{Solution 6 (super quick)}

Let $f(x) = \log(x + \log((x-1) + \log((x-2) + \ldots + \log 2 \ldots )))$. From the answer choices, we see that $3 < f(2013) < 4$. Since $f(x)$ grows very slowly, we can assume $3 < f(2012) < 4$. Therefore, $f(2013) = \log(2013 + f(2012)) \in (\log 2016, \log 2017) \implies \boxed{\textbf{(A)}}$.

 



\textbf{Solution 7 (Not rigorous but it works)}

Start small from $\log(3+\log2)$. Because $\log2$ is close to $0$, the expression is close to $\log 3$. This continues for the 1 digit numbers. However, when we get to $\log(11+\log10)$, $\log10=1$ so the expression equals $\log12$. We can round down how much each new log contributes to the expression as an approximate answer. There are 90 2 digit numbers, 900 3 digit numbers, and 1014 4 digit numbers that increase the expression. $1000< 90 \cdot 1 + 900 \cdot 2 + 1014 \cdot 3< 10000$, so it increases between 3 and 4. The answer is $\boxed{\textbf{(A)}}$. 

 



\textbf{Solution 8 (Assumption)}

Consider $\log (2013+\log(2012))$ and disregard every other log. We do this because we assume that the logs from $2011$ onward cannot contribute more than $7987$ because of the nature of logs. We have $\log(2016)<\log (2013+\log(2012))<\log(2017)$ so the answer is $\boxed{A}.$

 


 \textbf{= Solution 9 (i just lost my dawg)} It seems that $\log (2012+x)$ is going to be less than four, since otherwise, it would imply that $x>7987$, and going down, that seems very impossible.

 Then, the answer is just $\boxed{\textbf{A} (\log{2016}, \log{2017})}$

 Note: This solution is only for people with unlimited aura. (And veryyy limited time)

 


	\item (\textit{AMC 12 12A 2013}) A palindrome is a nonnegative integer number that reads the same forwards and backwards when written in base 10 with no leading zeros. A 6-digit palindrome $n$ is chosen uniformly at random. What is the probability that $\frac{n}{11}$ is also a palindrome?


 $\textbf{(A)} \ \frac{8}{25} \qquad \textbf{(B)} \ \frac{33}{100} \qquad \textbf{(C)} \ \frac{7}{20} \qquad \textbf{(D)} \ \frac{9}{25} \qquad \textbf{(E)} \ \frac{11}{30}$

 \textbf{Answer: }E



\textbf{Solution 1}

By working backwards, we can multiply 5-digit palindromes $ABCBA$ by $11$, giving a 6-digit palindrome:

 $\overline{A (A+B) (B+C) (B+C) (A+B) A}$

 Note that if $A + B >= 10$ or $B + C >= 10$, then the symmetry will be broken by carried 1s

 Simply count the combinations of $(A, B, C)$ for which $A + B < 10$ and $B + C < 10$
$A = 1$ implies $9$ possible $B$ (0 through 8), for each of which there are $10, 9, 8, 7, 6, 5, 4, 3, 2$ possible C, respectively. There are $54$ valid palindromes when $A = 1$
$A = 2$ implies $8$ possible $B$ (0 through 7), for each of which there are $10, 9, 8, 7, 6, 5, 4, 3$ possible C, respectively. There are $52$ valid palindromes when $A = 2$

 Following this pattern, the total is

 $54 + 52 + 49 + 45 + 40 + 34 + 27 + 19 + 10 = 330$

 6-digit palindromes are of the form $XYZZYX$, and the first digit cannot be a zero, so there are $9 \cdot 10  \cdot 10 = 900$ combinations of $(X, Y, Z)$

 So, the probability is $\frac{330}{900} = \frac{11}{30}$

 \textbf{Note} You can more easily count the number of triples $(A, B, C)$ by noticing that there are $9 - B$ possible values for $A$ and $10 - B$ possible values for $C$ once $B$ is chosen. Summing over all $B$, the number is 
 \[9\cdot 10 + 8\cdot 9 + \ldots + 1\cdot 2 = 2\left(\binom{2}{2} + \binom{3}{2} + \ldots + \binom{10}{2}\right).\]By the hockey-stick identity, it is $2\binom{11}{3} = 330$.

 



\textbf{Solution 2 (using the answer choices)}

Let the palindrome be the form in the previous solution which is $XYZZYX$. It doesn't matter what $Z$ is because it only affects the middle digit. There are $90$ ways to pick $X$ and $Y$, and the only answer choice with denominator a factor of $90$ is $\boxed{\textbf{(E)} \ \frac{11}{30}}$.

 



\textbf{Solution 3 (extremely simple)}

A six-digit palindrome, that when divided by 11, has to, if another palindrome, equal a five-digit palindrome. It cannot be a four-digit palindrome, otherwise, it would be too small. We now have a five-digit palindrome $ABCBA$, where the digits could be the same. If we multiply by $11$ we are adding like this:

 \begin{verbatim}
ABCBA0
\end{verbatim}

+ ABCBA=

 \begin{verbatim}
A, A+B, B+C, B+C, A+B, A
\end{verbatim}

Commas separating the digits,We don't have to worry about A because it can be any digit, however, we have to do casework on B, as it affects all the digits.

 B=0 A=1-9 C=0-9,B=1 A=1-8 C=0-8,B=2 A=1-7 C=0-7,B=3 A=1-6 C=0-6,B=4 A=1-5 C=0-5,B=5 A=1-4 C=0-4,B=6 A=1-3 C=0-3,B=7 A=1-2 C=0-2,B=8 A=1   C=0-1,

 B=9, does not work because A's only option is equal to 0 which makes it an invalid number (four-digit)

 These solutions work as we ensure that all the digits do not carry into the following digits (e.g., The tens cannot carry into the hundreds or it would not be a palindrome).

 For probability, as Richard Rusczyk says the total possibilities over the total successful outcomes Possibilities--> which are A=1-9, B=0-9, and C=0-9 --> 9 x 10 x 10 = 900

 Successful Outcomes--> which are from the casework above --> 
 \[(9 \cdot 10) + (8 \cdot 9) + (7 \cdot 8) + (6 \cdot 7) + (5 \cdot 6) + (4 \cdot 5) + (3 \cdot 4) + (2 \cdot 3) + (1 \cdot 2)\] = $90 + 72 + 56 + 42 + 30 + 20 + 12 + 6 + 2 = 330$ 

 330/900 = $\boxed{\textbf{(E)} \ \frac{11}{30}}$



	\item (\textit{AMC 12 12A 2013}) $ABCD$ is a square of side length $\sqrt{3} + 1$. Point $P$ is on $\overline{AC}$ such that $AP = \sqrt{2}$. The square region bounded by $ABCD$ is rotated $90^{\circ}$ counterclockwise with center $P$, sweeping out a region whose area is $\frac{1}{c} (a \pi + b)$, where $a$, $b$, and $c$ are positive integers and $\text{gcd}(a,b,c) = 1$. What is $a + b + c$?


 $\textbf{(A)} \ 15 \qquad \textbf{(B)} \ 17 \qquad \textbf{(C)} \ 19 \qquad \textbf{(D)} \ 21 \qquad \textbf{(E)} \ 23$

 \textbf{Answer: }C



\textbf{Solution}

We first note that diagonal $\overline{AC}$ is of length $\sqrt{6} + \sqrt{2}$. It must be that $\overline{AP}$ divides the diagonal into two segments in the ratio $\sqrt{3}$ to $1$. It is not difficult to visualize that when the square is rotated, the initial and final squares overlap in a rectangular region of dimensions $2\sqrt{3}$ by $\sqrt{3} + 1$. The area of the overall region (of the initial and final squares) is therefore twice the area of the original square minus the overlap, or $2 (\sqrt{3} + 1)^2 - 2 (\sqrt{3} + 1) = 2 (4 + 2 \sqrt{3}) - 2 \sqrt{3} - 2 = 6 + 2 \sqrt{3}$.

 
The area also includes $4$ circular segments. Two are quarter-circles centered at $P$ of radii $\sqrt{2}$ (the segment bounded by $\overline{PA}$ and $\overline{PA'}$) and $\sqrt{6}$ (that bounded by $\overline{PC}$ and $\overline{PC'}$). Assuming $A$ is the bottom-left vertex and $B$ is the bottom-right one, it is clear that the third segment is formed as $B$ swings out to the right of the original square [recall that the square is rotated counterclockwise], while the fourth is formed when $D$ overshoots the final square's left edge. To find these areas, consider the perpendicular from $P$ to $\overline{BC}$. Call the point of intersection $E$. From the previous paragraph, it is clear that $PE = \sqrt{3}$ and $BE = 1$. This means $PB = 2$, and $B$ swings back inside edge $\overline{BC}$ at a point $1$ unit above $E$ (since it left the edge $1$ unit below). The triangle of the circular sector is therefore an equilateral triangle of side length $2$, and so the angle of the segment is $60^{\circ}$. Imagining the process in reverse, it is clear that the situation is the same with point $D$.

 
The area of the segments can be found by subtracting the area of the triangle from that of the sector; it follows that the two quarter-segments have areas $\frac{1}{4} \pi (\sqrt{2})^2 - \frac{1}{2} \sqrt{2} \sqrt{2} = \frac{\pi}{2} - 1$ and $\frac{1}{4} \pi (\sqrt{6})^2 - \frac{1}{2} \sqrt{6} \sqrt{6} = \frac{3 \pi}{2} - 3$. The other two segments both have area $\frac{1}{6} \pi (2)^2 - \frac{(2)^2 \sqrt{3}}{4} = \frac{2 \pi}{3} - \sqrt{3}$.

 
The total area is therefore 
 \[(6 + 2 \sqrt{3}) + (\frac{\pi}{2} - 1) + (\frac{3 \pi}{2} - 3) + 2 (\frac{2 \pi}{3} - \sqrt{3})\] 
 \[= 2 + 2 \sqrt{3} + 2 \pi + \frac{4 \pi}{3} - 2 \sqrt{3}\] 
 \[= \frac{10 \pi}{3} + 2\] 
 \[= \frac{1}{3} (10 \pi + 6)\]

 
Since $a = 10$, $b = 6$, and $c = 3$, the answer is $a + b + c = 10 + 6 + 3 = \boxed{\textbf{(C)} \ 19}$.

 


	\item (\textit{AMC 12 12A 2013}) Three distinct segments are chosen at random among the segments whose end-points are the vertices of a regular $12$-gon. What is the probability that the lengths of these three segments are the three side lengths of a triangle with positive area?


 $\textbf{(A)} \ \frac{553}{715} \qquad \textbf{(B)} \ \frac{443}{572} \qquad \textbf{(C)} \ \frac{111}{143} \qquad \textbf{(D)} \ \frac{81}{104} \qquad \textbf{(E)} \ \frac{223}{286}$

 \textbf{Answer: }E



\textbf{Solution}

Suppose $p$ is the answer. We calculate $1-p$.

 Assume that the circumradius of the 12-gon is $1$, and the 6 different lengths are $a_1$, $a_2$, $\cdots$, $a_6$, in increasing order. Then

 $a_k = 2\sin ( \frac{k\pi}{12} )$. 

 So $a_1=(\sqrt{6}-\sqrt{2})/2 \approx 0.5$, 

 $a_2=1$,

 $a_3=\sqrt{2}\approx 1.4$, 

 $a_4=\sqrt{3}\approx 1.7$,

 $a_5=(\sqrt{6}+\sqrt{2})/2 = a_1 + a_3$,

 $a_6 = 2$.

 


 Now, Consider the following inequalities:

 $a_3>2a_1 > a_2$
$a_4> a_1 + a_2>a_3$
$a_4<a_1 + a_3=a_5$
$a_1 + a_4 > a_6$
$2a_2 = 2 = a_6$. Thus any two segments with at least one them longer than $a_2$ have a sum greater than $a_6$.

 Therefore, all triples (in increasing order) that can't be the side lengths of a triangle are the following. Note that x-y-z means $(a_x, a_y, a_z)$:

 \begin{verbatim}
1-1-3, 1-1-4, 1-1-5, 1-1-6,1-2-4, 1-2-5, 1-2-6,1-3-5, 1-3-6,2-2-6
\end{verbatim}

Note that there are $12$ segments of each length of $a_1$, $a_2$, $\cdots$, $a_5$, respectively, and $6$ segments of length $a_6$. There are $66$ segments in total.

 In the above list there are $3$ triples of the type a-a-b without \textit{6}, $2$ triples of a-a-6 where a is not \textit{6}, $3$ triples of a-b-c without \textit{6}, and $2$ triples of a-b-6 where a, b are not \textit{6}. So,

 
 \[1-p = \frac{1}{66\cdot 65\cdot 64} ( 3\cdot 3 \cdot 12\cdot 11\cdot 12 + 2\cdot 3 \cdot 12\cdot 11\cdot 6 + 3\cdot 6\cdot 12^3 + 2\cdot 6 \cdot 12^2 \cdot 6)\]

 \[= \frac{1}{66\cdot 65\cdot 64} (12^2 (99+33) + 12^3(18+6)) = \frac{1}{66\cdot 65\cdot 64} (12^3 \cdot 35) = \frac{63}{286}\]

 So $p = 223/286$.

 


	\item (\textit{AMC 12 12A 2013}) Let $f : \mathbb{C} \to \mathbb{C}$ be defined by $f(z) = z^2 + iz + 1$. How many complex numbers $z$ are there such that $\text{Im}(z) > 0$ and both the real and the imaginary parts of $f(z)$ are integers with absolute value at most $10$?


 $\textbf{(A)} \ 399 \qquad \textbf{(B)} \ 401 \qquad \textbf{(C)} \ 413 \qquad \textbf{(D)} \ 431 \qquad \textbf{(E)} \ 441$

 \textbf{Answer: }A



\textbf{Solution 1}

Suppose $f(z)=z^2+iz+1=c=a+bi$. We look for $z$ with $\operatorname{Im}(z)>0$ such that $a,b$ are integers where $|a|, |b|\leq 10$.

 First, use the quadratic formula:

 $z = \frac{1}{2} (-i \pm \sqrt{-1-4(1-c)}) = -\frac{i}{2} \pm \sqrt{ -\frac{5}{4} + c }$

 Generally, consider the imaginary part of a radical of a complex number: $\sqrt{u}$, where $u = v+wi = r e^{i\theta}$.

 $\operatorname{Im}(\sqrt{u}) = \operatorname{Im}(\pm \sqrt{r} e^{i\theta/2}) = \pm \sqrt{r} \sin(\theta/2) = \pm \sqrt{r}\sqrt{\frac{1-\cos\theta}{2}} = \pm \sqrt{\frac{r-v}{2}}$.

 Now let $u= -5/4 + c$, then $v = -5/4 + a$, $w=b$, $r=\sqrt{v^2 + w^2}$.

 Note that $\operatorname{Im}(z)>0$ if and only if $\pm \sqrt{\frac{r-v}{2}}>\frac{1}{2}$. The latter is true only when we take the positive sign, and that $r-v > 1/2$,

 or $v^2 + w^2 > (1/2 + v)^2 = 1/4 + v + v^2$, $w^2 > 1/4 + v$, or $b^2 > a-1$.

 In other words, when $b^2 > a-1$, the equation $f(z)=a+bi$ has unique solution $z$ in the region $\operatorname{Im}(z)>0$; and when $b^2 \leq a-1$ there is no solution. Therefore the number of desired solution $z$ is the same as the number of ordered pairs $(a,b)$ such that integers $|a|, |b|\leq 10$, and that $b^2 \geq a$.

 When $a\leq 0$, there is no restriction on $b$ so there are $11\cdot 21 = 231$ pairs;

 when $a > 0$, there are $2(1+4+9+10+10+10+10+10+10+10)=2(84)=168$ pairs.

 So there are $231+168=\boxed{399}$ in total.

 



\textbf{Solution 2 (motivated by coordinate geometry)}

We consider the function $f(z)$ as a mapping from the 2-D complex plane onto itself. We complete the square of $f(z)=z^2+iz+1=(z+\frac{i}{2})^2+\frac{5}{4}$.

 Now, we must decide the range of $f(z)$ based on the domain of $z$, $\operatorname{Im}(z)>0$. To do this, we are interested in mapping the boundary line $\operatorname{Im}(z)=0$. To make the mapping simpler, let $f(z)=g(z)+\frac{5}{4}$, or $g(z)=(z+\frac{i}{2})^2$.

 We intend to map of the line $\operatorname{Im}(z)=0$ using the function $g(z)$. This transformation is equivalent to the polar equation $r=(\frac{1}{2}\csc(\frac{\theta}{2}))^2$. Using polar and trig identities, we can restate this equation as the rectangular form of a parabola,

 $x=y^2-\frac{1}{4}$,

 where $x=\operatorname{Re}(z)$ and $y=\operatorname{Im}(z)$. So, we conclude that $f(z)$ maps the line $\operatorname{Im}(z)=0$ to the parabola 

 $x=y^2-\frac{1}{4}+\frac{5}{4}=y^2+1$.

 A quick check reveals that the range of $f(z)$ is to the left of the parabola, meaning that any point on or to the right of parabola cannot be reached.

 Since the problem requires $|\operatorname{Re}(z)|$ and $|\operatorname{Im}(z)|$ to both be integers and at most 10, all that remains is counting all points with integer coordinates in the range of $f(z), \operatorname{Im}(z)>0$. To do this, we employ complementary counting.

 The points of interest are $|\operatorname{Re}(z)|\leq 10$ and $|\operatorname{Im}(z)|\leq 10$, resulting in a total of $441$ points. For lattice points on or to the right of the parabola, there are $10$ points for $x=0$, $9$ points for $x=\pm 1$, $6$ points for $x=\pm 2$, and $1$ point for $x=\pm 3$. Summing it all together, our answer is $441-(10+2*9+2*6+2*1)=\boxed{399}$.

 



\textbf{Solution 3 (Pure Algebraic Manipulation, by Junfei)}

Let $z=a+bi$, where $a$ and $b$ are real numbers with $b>0$.

 
$z^2+iz+1$
$=(a+bi)^2+i(a+bi)+1$
$=a^2+2abi-b^2+ai-b+1$
$=(a^2-b^2-b+1)+(a+2ab)i$

 
The conditions are:

 $a^2-b^2-b+1$ is an integer with absolute value no greater than $10$, and

 $a+2ab$ is an integer with absolute value no greater than $10$.

 
Let $a^2-b^2-b+1=M$ and $a+2ab=N$.

 
$a(2b+1)=N$
$a=\frac{N}{2b+1}$ This fraction always exists because $2b+1>1$. This is the necessary and sufficient condition for the second equation.

 
Substitute $a=\frac{N}{2b+1}$ into the first equation.

 $(\frac{N}{2b+1})^2-b^2-b+1=M$
$\frac{4N^2}{(2b+1)^2}-4b^2-4b+4=4M$
$\frac{4N^2}{(2b+1)^2}-4b^2-4b-1=4M-5$
$\frac{4N^2}{(2b+1)^2}-(2b+1)^2=4M-5$

 
Let $(2b+1)^2=g$. Since it is given that $b>0$, it is necessary that $g>1$. If $g\leq{1}$, $b$ is either unreal or negative. If $g>1$, there exists a unique positive value of $b$ that satisfies $(2b+1)^2=g$.

 
$\frac{4N^2}{g}-g=4M-5$
$4N^2-g^2=(4M-5)g$
$g^2+(4M-5)g-4N^2=0$
$g=\frac{-(4M-5)\pm\sqrt{(4M-5)^2+16N^2}}{2}$ This is the necessary and sufficient condition for the first equation. 

 
Both equations are satisfied when we have $g=\frac{-(4M-5)\pm\sqrt{(4M-5)^2+16N^2}}{2}$ and $a=\frac{N}{2b+1}$, but the condition $b>0$, or $g>1$, has to be satisfied as well.

 
For $g$ to be a positive real number, we can only take the plus sign from $g=\frac{-(4M-5)\pm\sqrt{(4M-5)^2+16N^2}}{2}$. Taking the minus sign makes $g$ a non-positive number, contradicting the requirement that $b>0$.

 $g=\frac{-(4M-5)+\sqrt{(4M-5)^2+16N^2}}{2}>1$
$-(4M-5)+\sqrt{(4M-5)^2+16N^2}>2$
$\sqrt{(4M-5)^2+16N^2}>4M-3$

 This is always satisfied when $4M-3$ is negative. When $4M-3\geq 0$:

 $(4M-5)^2+16N^2>(4M-3)^2$
$-40M+25+16N^2>-24M+9$
$16N^2>16M-16$
$N^2>M-1$

 
This is the necessary and sufficient condition for $g>1$, or $b>0$.

 
For $M\leq{0}$, $N^2>M-1$ is always true.

 For $M=1$, $N$ has $20$ choices.

 For $M=2,3,4$, $N$ has $18$ choices.

 For $M=5,6,7,8,9$, $N$ has $16$ choices.

 For $M=10$, $N$ has $14$ choices.

 
The number of outputs of $f(z)$ is $11\times{21}+20+3\times{18}+5\times{16}+14=399$.

 
Recall the equations in which $g$, $b$ and $a$ are expressed using $M$ and $N$. $g=\frac{-(4M-5)+\sqrt{(4M-5)^2+16N^2}}{2}$, $b=\frac{\sqrt{g}-1}{2}$, and $a=\frac{N}{2b+1}$.

 
For each pair $(M,N)$ with $N^2>M-1$, a single value of $b>0$ is generated. For a given $N$ and a given $b>0$, a single value of $a$ is generated. In a nutshell, for each given pair $(M,N)$ with $N^2>M-1$, a single pair of $(m,n)$ is generated.

 
Furthermore, for each pair of $(M,N)$ without $N^2>M-1$, we either have an unreal $b$ or a non-positive $b$.

 
Since each of the $399$ pairs of $(M,N)$ results in exactly one pair of $(a,b)$, or exactly one value of $z=a+bi$, the total number of values of $z$ is also $399$. This completes the proof.

 


\end{enumerate}